\section{Related work}
\label{sec:related}

\paragraph{Hypergraph and n-ary data stores.} HyperGraphDB~\cite{hgdb} is the closest system and our baseline. Its atoms are identified by UUID handles, a link is an ordered array of target handles, and an index maps every atom to its incidence set; storage is delegated to Berkeley DB Java Edition. A link is immutable, so changing a membership replaces the link, and the incidence index has to be updated for every member. OpenCog's AtomSpace~\cite{opencog} stores typed atoms and links in memory for reasoning workloads. TypeDB~\cite{typeql} gives $n$-ary relations and their role players a type system and a query language, and stores them in a key-value engine. Property-graph databases and RDF stores represent $n$-ary relations by reification~\cite{angles}, and Wikidata attaches qualifiers to statements for the same purpose~\cite{wikidata}. HStore differs from all of them in that a hyperedge's membership is a storage structure of its own, versioned, counted and summarised like every other tree in the engine.

\paragraph{Hypergraph processing.} Hygra~\cite{hygra}, MESH~\cite{mesh} and HyperX~\cite{hyperx} run analytics over hypergraphs held in memory, on one machine or a cluster. They assume a static input and optimise traversal and message passing. HStore is a storage engine for data that changes, and the analytic operators in its database layer read from its snapshots.

\paragraph{Copy-on-write and shadow paging.} Shadow paging goes back to Lorie~\cite{lorie}. Rodeh showed how to make B-trees copy-on-write with reference counts so that snapshots and clones are cheap~\cite{rodeh-shadowing}, and Btrfs builds a file system on that design~\cite{btrfs}. LMDB~\cite{lmdb} uses a copy-on-write B+tree with a single writer and readers that pin a root, much like HStore's generations. These systems write the changed path at every commit. HStore defers those writes to checkpoints and logs keys instead, which is closer to how a log-structured store absorbs small writes, while keeping the property that a published version never changes. Making data structures persistent in the functional sense~\cite{persistent} and the owner tokens that let one transaction change its own fresh nodes in place follow the same idea at the level of single objects.

\paragraph{Logging, recovery and log-structured storage.} ARIES~\cite{aries} logs page-level changes and redoes them in place. HStore never changes a page in place, so it logs logical key changes and replays them onto the last checkpoint's trees, and its only physical log records are page references for the rare commits that write pages. The rule that recovery stops at the first commit whose pages are missing addresses the reordering of writes by the file system that crash-consistency studies have catalogued~\cite{pillai,torturing}. Log-structured file systems~\cite{lfs} and LSM-trees~\cite{lsm} append all writes and reclaim space by cleaning or merging; Berkeley DB Java Edition, under HyperGraphDB, is a log-structured B+tree whose cleaner copies live records forward. HStore's compaction also copies live data forward, but because trees refer to pages by number through a directory, it copies pages without rewriting their parents. The Bw-tree~\cite{bwtree} uses a mapping table from logical page ids to physical addresses for a related reason, to install delta updates and relocate pages without changing references. LeanStore~\cite{leanstore} replaces the buffer pool's hash table with pointer swizzling; HStore's node cache stays a separate, budgeted table because its nodes are immutable and shared by every reader.

\paragraph{Order maintenance and set algebra.} Ordered hyperedges use the order-maintenance labelling of Dietz and Sleator and its later simplifications~\cite{dietz-sleator,bender-order}, as ORDPATH labels do for XML~\cite{ordpath}. Intersections use a leapfrog join~\cite{leapfrog}. Additive multiset hashes~\cite{multiset-hash} and polynomial fingerprints~\cite{karp-rabin} give each tree a content fingerprint that is independent of its shape.
